% Extracto íntegro por residencia del propietario científico definitivo de masas.
\section{Domaine interne et problème de réalisation}

La Loi générale des masses ne part pas d'une table de noms ni d'une liste de
valeurs expérimentales. Son domaine interne est la chaîne finie
\[
 \mathcal R_{324}\longrightarrow \mathcal C_{81}
 \longrightarrow \mathcal F_{56}\longrightarrow \mathcal M_{13},
\]
où \(\mathcal R_{324}\) contient les \(4\) orientations de chacune des
\(81\) cellules, \(\mathcal F_{56}\) réunit des histoires de routes équivalentes et
\(\mathcal M_{13}\) conserve les multisections opératorielles. Le secteur nul
bifurque avant l'application du caractère positif. Par conséquent, ni une table
historique abrégée ni une taxonomie de noms physiques ne peuvent remplacer ce
domaine.

Soit \(\mathcal P\) l'ensemble des états physiques individualisés. Un état
élémentaire est représenté par le descripteur
\[
 X=(\mathfrak c,\rho,g,f,q,\epsilon,\varsigma,\mathfrak a),
\]
qui conserve la classe de représentation \(\mathfrak c\), la représentation interne
\(\rho\), la génération \(g\), la saveur \(f\), la charge \(q\), l'orientation de feuillet
\(\epsilon\), le secteur statistique \(\varsigma\) et les données de couplage
\(\mathfrak a\). Un état composé ajoute des constituants ordonnés, un arbre de
composition, une frontière commune et des nombres quantiques collectifs. Une résonance
ajoute les canaux ouverts et la prescription causale de la résolvante.

La réalisation physique doit répondre à trois questions distinctes :

\begin{enumerate}
\item quelle fibre, orbite ou multisection de l'atlas correspond au descripteur ;
\item quel opérateur positif publie la masse sur cette fibre ;
\item quel couplage produit un scalaire, un multiplet, une distribution ou
      un pôle complexe.
\end{enumerate}

Confondre ces questions introduit un sélecteur rétrospectif. La réponse
correcte s'obtient par une relation naturelle, un opérateur de fibre et une
compression spectrale, dans cet ordre.

\section{Produit fibré entre descripteurs et routes}

Soit \(\mathcal D\) l'espace fini des données internes déjà construites avant la
confrontation : représentation, orientation, mémoire, frontière, charge, saveur,
couleur, génération et secteur statistique. Notons
\[
 d_{\mathcal P}:\mathcal P\longrightarrow\mathcal D,
 \qquad
 d_{\mathcal R}:\mathcal R_{324}\longrightarrow\mathcal D
\]
les applications de descripteur physique et de descripteur généalogique. On définit l'objet
de réalisation par le produit fibré
\[
 \mathcal J
 =\mathcal P\times_{\mathcal D}\mathcal R_{324}
 =\{(X,\gamma):d_{\mathcal P}(X)=d_{\mathcal R}(\gamma)\}.
\]
La fibre d'une espèce est
\[
\mathscr S(X)=\{\gamma\in\mathcal R_{324}:(X,\gamma)\in\mathcal J\}.
\]
Un descripteur (X) est dit réalisable par l'atlas lorsque
\[
 d_{\mathcal P}(X)\in\operatorname{im}d_{\mathcal R}.
\]
Cette condition équivaut à \(\mathscr S(X)\ne\varnothing\). Elle ne se déduit pas
de la simple écriture du produit fibré : elle constitue la condition exacte d'
existence qui doit être vérifiée pour chaque descripteur physique.
La valeur de \(\mathscr S\) peut être une route, une orbite ou une multisection.
On ne force pas une section ponctuelle lorsque la symétrie interne exige de conserver
une multiplicité.

\begin{definition}[Relation de réalisation admissible]
La relation \(\mathscr S\) est admissible lorsqu'elle satisfait simultanément :
\begin{enumerate}
\item elle dépend uniquement de données construites avant l'ouverture de la table extérieure ;
\item elle commute avec la conjugaison particule--antiparticule ;
\item elle est équivariante sous l'inversion cellulaire et les actions internes de
      jauge ;
\item elle est compatible avec le raffinement, le retour nonadique et la lecture
      dodécaphasique ;
\item elle conserve toute multiplicité non éliminée par un couplage physique ;
\item elle reste invariante lorsqu'on permute ou remplace les valeurs extérieures de
      masse et de largeur.
\end{enumerate}
\end{definition}

\begin{theorem}[Réalisation prospective par produit fibré]
Si \(d_{\mathcal P}\) et \(d_{\mathcal R}\) sont équivariantes par rapport à une
symétrie \(G\) et ne dépendent pas de grandeurs extérieures, alors \(\mathcal J\)
est un \(G\)-sous-objet de \(\mathcal P\times\mathcal R_{324}\). Pour chaque
descripteur réalisable \(X\), la fibre non vide \(\mathscr S(X)\) est
\(G_X\)-stable et l'assignation \(X\mapsto\mathscr S(X)\) est indépendante du
résultat auquel elle sera ultérieurement confrontée.
\end{theorem}

\begin{proof}
Pour \((X,\gamma)\in\mathcal J\), on a
\(d_{\mathcal P}(X)=d_{\mathcal R}(\gamma)\). L'équivariance donne
\[
 d_{\mathcal P}(gX)=g\,d_{\mathcal P}(X)
 =g\,d_{\mathcal R}(\gamma)=d_{\mathcal R}(g\gamma),
\]
et, par conséquent, \((gX,g\gamma)\in\mathcal J\). Si \(g\in G_X\), la
seconde composante reste dans \(\mathscr S(X)\). L'indépendance des
valeurs extérieures est immédiate, car aucune d'elles n'appartient au domaine
des deux applications du produit fibré.
\end{proof}

L'électron constitue l'exception centrale. L'inversion cellulaire possède un
unique point fixe, la cellule \((5,5)\), et sa fibre physique est donc de rang
un. Les autres multisections stables n'admettent pas, en général, de sélecteur
ponctuel équivariant. Cette différence n'est pas une lacune de l'atlas : c'est la
distinction mathématique entre une section centrale et une espèce réalisée comme
orbite ou multisection.

\section{Classification spectrale intrinsèque des espèces}

Soit \(\mathcal A_{\rm rut}\) l'algèbre commutative finie engendrée par les
observables internes de route : signature \(\mathbb Z^5\), caractère, chronologie,
\(K_D\), \(K_\Omega\), orientation, mémoire et lecteurs compatibles. Dans une
fibre \(F\), on définit
\[
 \gamma\sim_{\rm sp}\eta
 \quad\Longleftrightarrow\quad
 a(\gamma)=a(\eta)\quad\text{pour tout }a\in\mathcal A_{\rm rut}.
\]
Pour chaque classe \(\lambda\in F/\!\sim_{\rm sp}\),
\[
 P_\lambda=\sum_{\gamma\in\lambda}
 |\gamma\rangle\langle\gamma|
\]
est un projecteur spectral ; les \(P_\lambda\) sont orthogonaux et
\(\sum_\lambda P_\lambda=I_F\).

\begin{theorem}[Individualisation interne maximale]
Tout classificateur construit exclusivement avec des observables de
\(\mathcal A_{\rm rut}\) se factorise de manière unique par le quotient
\(F/\!\sim_{\rm sp}\). Par conséquent, deux descripteurs de signatures complètes
distinctes appartiennent à des secteurs distinguables ; deux descripteurs qui coïncident
pour tous les générateurs de l'algèbre ne peuvent pas être séparés par une règle
construite avec ces mêmes observables.
\end{theorem}

\begin{proof}
Les caractères de l'algèbre commutative finie séparent exactement son
spectre conjoint. Un classificateur fonctionnel de ses générateurs est constant
sur chaque fibre de l'application spectrale conjointe et se factorise donc par le
quotient. Si deux routes ont le même caractère sur toute l'algèbre, aucun
élément de celle-ci ne peut les séparer.
\end{proof}

Ce théorème clôt la limite logique de la classification. Lorsqu'une
représentation physique exige de distinguer deux états appartenant à une même
classe spectrale, elle doit exhiber un observable supplémentaire ou un couplage qui
étend \(\mathcal A_{\rm rut}\). L'absence d'un tel observable n'autorise pas à
utiliser la masse extérieure comme étiquette de sélection.

Dans le secteur leptonique chargé, les descripteurs distinguent la fibre centrale,
la multisection de profondeur \(G_2\) et la multisection de profondeur
\(G_4\). La première produit l'unique ligne électronique ; les deux autres
produisent les opérateurs associés aux réalisations muonique et tauonique. Si
le couplage préserve plus d'un bloc, la prédiction correcte est son
spectre ; l'étiquette physique n'autorise pas à choisir l'une de ses valeurs propres par
proximité numérique.


